\chapter{Reverse-time diffusion and generation}
\label{ch:reverse-sde-generation}

A diffusion model generates by running the channel of \cref{ch:ou-channel-tweedie}
backwards in time, and the only ingredient the backward equation needs that the
forward one does not is the score. This chapter derives the reverse-time equation
(Anderson's theorem) from the Fokker--Planck equation, explains why an error in the
score is exactly what an error in the generated law is made of, writes down the
Euler--Maruyama scheme the code uses, and then describes the paper's generation
experiment E05: its same-Brownian-path design, its arms, and every metric by which the
samples are judged, including the transition-law metric (Hellinger distance) used in
the learning section.

\section{Anderson's reverse-time equation}
\label{sec:rev:anderson}

Consider a forward diffusion on $[0,T]$,
\begin{equation}
  \dd\vx_t=f(\vx_t,t)\,\dd t+g(t)\,\dd\vect{W}_t,\qquad \vx_0\sim\Pz,
  \label{eq:rev:forward}
\end{equation}
with smooth positive marginal densities $p_t$. The evolution of $p_t$ is the
\term{Fokker--Planck equation}
\begin{equation}
  \partial_tp_t=-\nabla\cdot\bigl(f\,p_t\bigr)+\tfrac12g^2\,\Delta p_t ,
  \label{eq:rev:fp}
\end{equation}
where $\nabla\cdot$ is the divergence and $\Delta=\nabla\cdot\nabla$ the Laplacian
(not to be confused with $\Dt$). The reader who has not met \cref{eq:rev:fp} may read
it as a conservation law: the drift transports probability with velocity $f$, and the
noise spreads it by diffusion with coefficient $\tfrac12g^2$.

\begin{theorem}[Time reversal]
\label{thm:rev:anderson}
Let $\tau=T-t$ and let $\bar{\vx}_\tau$ solve
\begin{equation}
  \dd\bar{\vx}_\tau=\bigl[-f(\bar{\vx}_\tau,T-\tau)+g(T-\tau)^2\,\nabla\log p_{T-\tau}(\bar{\vx}_\tau)\bigr]\dd\tau
   +g(T-\tau)\,\dd\bar{\vect{W}}_\tau,\qquad \bar{\vx}_0\sim p_T,
  \label{eq:rev:reverse}
\end{equation}
with $\bar{\vect{W}}$ a Brownian motion. Then $\bar{\vx}_\tau$ has the law $p_{T-\tau}$
for every $\tau\in[0,T]$; in particular $\bar{\vx}_T\sim\Pz$. Anderson (1982) proved
more: the process $t\mapsto\vx_{T-t}$ itself is a solution of \cref{eq:rev:reverse}.
\end{theorem}

\begin{proof}[Proof of the marginal statement]
Scalar case, and write $p=p_t(x)$. Since $\tfrac12g^2\partial_x^2p=\partial_x\bigl(g^2\partial_xp\bigr)-\tfrac12g^2\partial_x^2p$
and $\partial_xp=p\,\partial_x\log p$, the Fokker--Planck equation can be rewritten as
\[
  \partial_tp=-\partial_x\Bigl[\bigl(f-g^2\partial_x\log p\bigr)p\Bigr]-\tfrac12g^2\,\partial_x^2p .
\]
Now set $q(x,\tau)=p_{T-\tau}(x)$, so that $\partial_\tau q=-\partial_tp$:
\[
  \partial_\tau q=-\partial_x\Bigl[\bigl(-f+g^2\partial_x\log q\bigr)q\Bigr]+\tfrac12g^2\,\partial_x^2q .
\]
This is the Fokker--Planck equation of the diffusion with drift $-f+g^2\partial_x\log q$
and diffusion coefficient $g$, run in the time $\tau$, and $q(\cdot,0)=p_T$. The
solution of \cref{eq:rev:reverse} has marginals obeying the same equation with the
same initial condition, hence the same marginals, by uniqueness for this parabolic
equation. In $\nsites$ dimensions replace $\partial_x$ by $\nabla\cdot$ and
$\partial_x^2$ by $\Delta$ line by line.
\end{proof}

A useful by-product of the same rewriting is the \term{probability-flow ODE}:
$\tfrac12g^2\Delta p=\nabla\cdot(\tfrac12g^2p\nabla\log p)$ turns \cref{eq:rev:fp} into
$\partial_tp=-\nabla\cdot[(f-\tfrac12g^2\nabla\log p)p]$, the transport equation of the
deterministic flow $\dot{\vx}=f-\tfrac12g^2\nabla\log p_t$, which has the same
marginals as the SDE and is implemented in \code{probability\_flow\_ode} of
\code{src/markov\_score/reverse.py}.

\begin{corollary}[The channel reversed]
\label{cor:rev:ou}
For the channel $f(\vx,t)=-\vx$, $g=\sqrt2$, the reverse equation in the original time
$t$ running from $T$ down to $0$ is
\begin{equation}
  \dd\vx=\bigl[-\vx-2\,\score(\vx,t)\bigr]\dd t+\sqrt2\,\dd\bar{\vect{W}}_t\qquad(\dd t<0),
  \label{eq:rev:ou}
\end{equation}
that is, in the forward-running time $\tau$, $\dd\bar{\vx}=[\bar{\vx}+2\score(\bar{\vx},T-\tau)]\dd\tau+\sqrt2\,\dd\bar{\vect{W}}_\tau$.
\end{corollary}

The drift has two parts. The term $\bar{\vx}$ undoes the contraction of the forward
drift; the term $2\score$ pushes towards regions where $\Pt$ is large. At high noise,
where $\score\approx-\vx$, the two nearly cancel and the reverse process is again an
Ornstein--Uhlenbeck process (\cref{ex:rev:gaussian}); at low noise the score carries
all the structure of $\Pz$.

\section{Why the score is exactly what drives generation}
\label{sec:rev:why}

Everything in \cref{eq:rev:ou} is known except $\score$: the drift $-\vx$ and the
coefficient $\sqrt2$ are the channel's. A generator that replaces $\score$ by an
estimate $\widehat s$ runs a different diffusion and produces a different law, and the
discrepancy is controlled by the score error alone.

\begin{proposition}[Score error controls the generated law]
\label{prop:rev:girsanov}
Let $\Prob$ be the law of the paths of \cref{eq:rev:ou} on $[\tmin,T]$ started at
$\vx_T\sim P_T$, and $\mathbb{Q}$ the law of the paths of the same equation with
$\widehat s$ in place of $\score$ and started at $\gamma=\normal{0}{\Id}$. Then
\begin{equation}
  \KL\bigl(P_{\tmin}\,\big\|\,\mathbb{Q}_{\tmin}\bigr)
  \le\KL\bigl(P_T\,\big\|\,\gamma\bigr)+\int_{\tmin}^{T}\E_{\Pt}\bigl|\widehat s(\vx_t,t)-\score(\vx_t,t)\bigr|^2\dd t ,
  \label{eq:rev:kl}
\end{equation}
where $\mathbb{Q}_{\tmin}$ is the law of the generated sample at time $\tmin$.
\end{proposition}

\begin{proof}[Proof sketch]
Two diffusions with the same diffusion coefficient $g$ and drifts $\mu_1,\mu_2$ have
path laws whose Kullback--Leibler divergence is, by Girsanov's theorem,
$\tfrac12\int\E|\mu_1-\mu_2|^2/g^2\,\dd t$ plus the divergence of the initial laws.
Here $\mu_1-\mu_2=2(\score-\widehat s)$ and $g^2=2$, so the integrand is
$\tfrac12\cdot4|\score-\widehat s|^2/2=|\score-\widehat s|^2$, with the expectation
under the true path law, whose marginal at time $t$ is $\Pt$. The divergence of the
marginals at $\tmin$ is at most that of the path laws (data processing).
\end{proof}

This is Theorem 1 of Song, Durkan, Murray and Ermon (2021). It says that the
time-integrated mean squared score error under the true noised laws is the quantity a
generator pays for, in nats. The paper's estimand (\cref{def:sm:estimand}) is the
square root of that integrand summed over the twelve levels and divided by the
reference energy; the generation experiment E05 measures whether the ordering of
estimators by score error is the ordering by sample quality, which
\cref{eq:rev:kl} predicts but does not guarantee, because it is an upper bound.

\section{The Euler--Maruyama scheme and its grid}
\label{sec:rev:euler}

\begin{definition}[Euler--Maruyama for the reverse channel]
\label{def:rev:em}
Given a decreasing sequence of times $T=t_0>t_1>\dots>t_n=\tmin$ with steps
$h_k=t_k-t_{k+1}>0$, the scheme is
\begin{equation}
  \vx_{k+1}=\vx_k+h_k\bigl[\vx_k+2\,\widehat s(\vx_k,t_k)\bigr]+\sqrt{2h_k}\,\vect{\xi}_k,
  \qquad\vect{\xi}_k\sim\normal{0}{\Id}\ \text{i.i.d.},\qquad \vx_0\sim\normal{0}{\Id}.
  \label{eq:rev:euler}
\end{equation}
\end{definition}

This is \code{reverse\_sde} and \code{reverse\_sde\_with\_increments} in
\code{src/markov\_score/reverse.py}: the drift is evaluated at the left endpoint, and
the Brownian increment over the step is $\sqrt{h_k}\,\vect{\xi}_k$, scaled by the
$\sqrt2$ of the channel. Three choices need explaining.

\paragraph{The grid is geometric.} The code uses $t_k=\tmax\,r^{k}$ with
$r=(\tmin/\tmax)^{1/n}$ (\code{time\_grid}), so that $h_k=(1-r)\,t_k$: the step is a
fixed fraction of the current time. The reason is that the reverse drift can be stiff
at small $t$: by \cref{eq:ou:tweedie2}, $\nabla\score=-\Id/\Dt+(e^{-2t}/\Dt^2)\Cov(\va\mid\vx)$,
whose entries can be of size $1/\Dt\approx1/(2t)$ when the posterior is narrower than
the likelihood (a kinked or discrete prior), and an explicit scheme is only stable for
steps below the reciprocal of the stiffness. A step proportional to $t$ keeps
$h_k/\Delta_{t_k}\approx(1-r)/2$ fixed. With the paper's $\tmax=3$, $\tmin=0.02$ and
$n=400$ steps (\code{experiments/e05\_generation.py}), $\ln(\tmin/\tmax)=-5.0106$,
$r=e^{-0.012527}=0.98755$, the first step is $h_0=3\times0.01245=0.0374$ and the last
is $h_{399}=0.0203\times0.01245=0.00025$; the ratio $h_k/\Delta_{t_k}$ is about $0.006$
throughout.

\paragraph{The start is $\normal{0}{\Id}$, not $P_T$.} The true law at $T=3$ is
not available to sample from, and $\gamma$ is its limit (\cref{prop:ou:limit}). The
cost is the first term of \cref{eq:rev:kl}. For the covariance twin it is computable:
$\Sig_T=\Id+e^{-2T}(\Sig_0-\Id)$ and, to second order in $e^{-2T}$,
$\KL(\normal{0}{\Sig_T}\|\gamma)=\tfrac12[\tr\Sig_T-\nsites-\log\det\Sig_T]\approx\tfrac14e^{-4T}\Frob{\Sig_0-\Id}^2$.
For the frozen chain, $\Frob{\Sig_0-\Id}^2=2\sum_{k=1}^{31}(32-k)\corr^{2k}\approx148$
at $\corr=0.85$ (a geometric sum), so the start costs about $\tfrac14e^{-12}\cdot148\approx2\times10^{-4}$
nats: negligible next to the discretisation, and the reason $\tmax=3$ suffices.

\paragraph{The stop is $\tmin=0.02$, not $0$.} At $t=0$ the score of a law with a
density that is not smooth (the Laplace chain has a kink) is singular, and for a law
without a density it does not exist. The generated sample is therefore a sample at
$\tmin$ and is compared with a reference sample of the true law at $\tmin$, obtained
by drawing clean sequences and noising them once (\code{rng\_ref} in
\code{\_gen\_cell}); comparing with clean sequences would be a category error. The
alternative readout $\pmean{\va}=e^{t}(\vx+\Dt\score)$ of \cref{thm:ou:tweedie}
(\code{denoising\_readout}) returns a posterior mean, which is not a sample.

\paragraph{Order of accuracy.} For an additive noise, as here, Euler--Maruyama has
strong order one (the pathwise error is $O(h)$) and weak order one (errors of
expectations are $O(h)$); the discretisation bias is measured, not assumed, by the
step study below.

\section{The same-Brownian-path design of E05}
\label{sec:rev:design}

The experiment (\code{experiments/e05\_generation.py}) is organised in cells
$(\text{seed},\nseq)$. In each cell one initial sample $\vx_0$ of $2000$ sequences and
one array of Brownian increments for all $400$ steps are drawn once, and every arm is
integrated by \code{reverse\_sde\_with\_increments} on exactly those increments. The
arms differ only in the score function $\widehat s$:
\begin{itemize}
\item \emph{exact}: the tree-exact score under the true kernel (grid message passing;
  its distance to the true law is the sampler's own discretisation floor);
\item \emph{em--bp}: the mixture kernel fitted in E02a for the same seed and $\nseq$,
  at its validation-selected checkpoint, driven through the same message passing;
\item \emph{lmmse}: the covariance twin, $\widehat s=-\Sig_t^{-1}\vx$
  (\cref{cor:ou:gaussian});
\item \emph{isotropic}: $\widehat s=-\vx$, the score of $\gamma$;
\item \emph{net}: network checkpoints from E03's confirmatory part for the same cell,
  converted from posterior means by \cref{eq:ou:score-mean-identity}.
\end{itemize}
Because the noise is shared, the difference between two arms' samples is due to their
scores and not to Monte Carlo scatter, and the metrics can be compared within cell and
then paired across seeds, as everywhere else in the paper. The reference sample $R$ of
$4000$ sequences at $\tmin$ is drawn from its own stream. The \emph{step} part of the
experiment integrates the exact arm at $100$, $200$, $400$ and $800$ steps on
\term{nested Brownian paths}: the increments of the finest grid are drawn once and the
coarse grid's increment over an interval is the sum of the fine increments inside it
(\code{nested\_brownian\_path}, \code{coarsen\_increments}). Two resolutions are then
driven by the same Brownian path and the root mean square of their difference is a
\term{strong error}, a pure measure of discretisation. With independent noise at each
resolution the difference would be dominated by sampling scatter and the study would
show only robustness, which is the mistake the code's docstring records.

\section{The sample-quality metrics}
\label{sec:rev:metrics}

Let $G\in\Reals^{n_g\times\nsites}$ be the generated sample and $R\in\Reals^{n_r\times\nsites}$
the reference sample, both at $\tmin$. All definitions are those of
\code{src/markov\_score/gen\_metrics.py} (\code{compare}); the paper's generation
table prints the first, the second, the fourth and the sliced Wasserstein distance.

\paragraph{Residuals.} For a sample $Y$ the \term{autoregressive residuals} are
$r_{j,i}=y_{j,i+1}-\corr\,y_{j,i}$ with the true $\corr$, $i=1,\dots,\nsites-1$. At
$\tmin=0$ they would be the innovations $\varepsilon_{i+1}$; at $\tmin>0$ they are
$e^{-\tmin}\varepsilon_{i+1}+\sqrt{\Delta_{\tmin}}(z_{i+1}-\corr z_i)$.

\paragraph{Excess kurtosis.} For a pooled sample $\{u\}$ with central moments $m_2,m_4$,
$\widehat\kappa=m_4/m_2^2-3$. Applied to the residuals it measures the non-Gaussian
shape of the innovation that the sampler reproduces; a Gaussian-innovation sampler
(the twin) gives zero up to sampling error. Its standard error is a bootstrap over
sequences ($200$ resamples).

\begin{proposition}[Residual kurtosis at $\tmin$]
\label{prop:rev:kurt}
If the innovation has variance $\ivar$ and excess kurtosis $\kappa$, then the residual
of the true law at level $t$ has
\[
  \Var(r)=e^{-2t}\ivar+\Dt(1+\corr^2),\qquad
  \operatorname{cum}_4(r)=e^{-4t}\kappa\,\ivar^2,\qquad
  \text{excess kurtosis}=\frac{\kappa\,e^{-4t}\ivar^2}{\bigl(e^{-2t}\ivar+\Dt(1+\corr^2)\bigr)^2}.
\]
\end{proposition}

\begin{proof}
$r=e^{-t}\varepsilon+\sqrt{\Dt}(z'-\corr z)$ is a sum of independent terms; variances
add, and fourth cumulants add with the Gaussian terms contributing nothing.
\end{proof}

So the reference row of the paper's table does not show the innovation's excess
kurtosis $3$ but its attenuated and diluted value at $\tmin$; the worked example
computes it. The same holds for every arm, which is why the table is read as a
comparison with the reference row and not with $3$.

\paragraph{Lag fourth cumulant.} $\operatorname{cum}(r_i,r_i,r_{i+1},r_{i+1})=\E[r_i^2r_{i+1}^2]-\E[r_i^2]\E[r_{i+1}^2]-2\E[r_ir_{i+1}]^2$
for centred residuals, pooled over sites. For the true law it is exactly zero at every
$t$: $r_i$ and $r_{i+1}$ share only the Gaussian $z_{i+1}$, joint cumulants are additive
over independent components, and a Gaussian has no fourth cumulant. It is therefore a
check for spurious dependence between neighbouring residuals, which a sampler with a
wrong score can create; it is not a signature of the innovation's shape.

\paragraph{Marginal Wasserstein-1.} Pooling all entries of $G$ and of $R$ into two
one-dimensional samples with quantile functions $F_G^{-1}$, $F_R^{-1}$,
$W_1=\int_0^1|F_G^{-1}(u)-F_R^{-1}(u)|\dd u$, approximated by the mean over $1999$
equally spaced $u$. It compares the site marginals only.

\paragraph{Sliced Wasserstein-2.} For $256$ random unit vectors $\vect{\theta}_k$
(common seed), project both samples, $G\vect{\theta}_k$ and $R\vect{\theta}_k$, compute
the one-dimensional $W_2^2$ of the projections by quantile matching,
$W_2^2(\mu,\nu)=\int_0^1(F_\mu^{-1}-F_\nu^{-1})^2\dd u$, and report
$\bigl(\tfrac1{256}\sum_kW_2^2\bigr)^{1/2}$. It sees the joint law through its
one-dimensional projections, and is dominated by second-order structure: two laws with
the same covariance and different fourth cumulants project to nearly the same
one-dimensional laws, which is why the twin and the exact arm are close in this
column.

\paragraph{Energy distance.} $2\E|G-R|-\E|G-G'|-\E|R-R'|$ with independent copies,
estimated on subsamples of $1000$ (Sz\'ekely and Rizzo); it is zero if and only if the
laws coincide, and is again dominated by low-order structure in practice.

\paragraph{Covariance errors.} The sample covariance $\widehat{\Sig}$ of $G$ against a
target $\Sig$: the relative Frobenius error $\Frob{\widehat{\Sig}-\Sig}/\Frob{\Sig}$
and the worst absolute error of the mean diagonal at lags $0$ to $5$. The code passes
$\Sig_0$ as the target although the sample lives at $\tmin$, where the covariance is
$\Sig_{\tmin}=e^{-2\tmin}\Sig_0+\Delta_{\tmin}\Id$; the difference is
$\Delta_{\tmin}(\Sig_0-\Id)$, whose relative Frobenius size at $\tmin=0.02$ is about
$0.039\times\Frob{\Sig_0-\Id}/\Frob{\Sig_0}\approx0.035$ for the frozen chain, a floor
the exact arm cannot go below. The paper's table does not print this column.

\paragraph{Residual autocorrelation.} $\rho_k=\E[r_ir_{i+k}]/\E[r_i^2]$ for
$k=1,\dots,5$ on the pooled centred residuals; the reported number is the largest
absolute difference between the generated and the reference values. For the true law
at $\tmin$ the residuals are not white ($\rho_1=-\corr\Delta_{\tmin}/\Var(r)$, from the
shared $z_{i+1}$), which is why the reference values are subtracted rather than zero.

\paragraph{Histogram divergence of the residuals.} With $60$ bins between the $0.001$
and $0.999$ quantiles of the reference residuals, $p$ the reference histogram and $q$
the generated one with add-one smoothing, $\sum_{p_b>0}p_b\log(p_b/q_b)$: a
Kullback--Leibler divergence of the innovation shape, true against generated.

\paragraph{Log-likelihood under the true law.} For every sequence of $G$ and of $R$
the evidence $\log P_{\tmin}(\vx)$ under the true kernel is computed exactly for the
discretised model by the same message passing that gives the score (the accumulated
normalisation constants of the recursion; the function \code{log\_\allowbreak likelihood\_\allowbreak true}), and the
reported gap is
\begin{equation}
  \Delta\ell=\frac1{n_r}\sum_{R}\log P_{\tmin}(\vx)-\frac1{n_g}\sum_{G}\log P_{\tmin}(\vx).
  \label{eq:rev:llgap}
\end{equation}
Writing $P$ for the true law and $Q$ for the generated one, in expectation
$\Delta\ell=\E_P\log P-\E_Q\log P=\KL(Q\|P)+H(Q)-H(P)$ with $H$ the differential
entropy: a difference of cross-entropies, not a divergence. It is zero when $Q=P$,
positive when $Q$ is misplaced or too spread, and can be negative when $Q$ is more
concentrated than $P$ on the high-density region. The exact arm's value is the
discretisation floor of the sampler, since its only error is the scheme's; the
superseded four-seed run recorded a negative value for that arm, which is consistent
with the definition and is the reason the sign is worth stating.

\section{A transition-law metric: the Hellinger distance}
\label{sec:rev:hellinger}

The learning section of the paper reports that the transition law is recovered ``to
Hellinger distance'' a stated value at the lowest level. The metric is the
function \code{transition\_\allowbreak hellinger} of
\code{src/\allowbreak markov\_score/\allowbreak metrics.py}.

\begin{definition}[Hellinger distance]
\label{def:rev:hellinger}
For two densities $p,q$ on $\Reals$,
\[
  H^2(p,q)=\tfrac12\int\bigl(\sqrt{p}-\sqrt{q}\bigr)^2\dd u=1-\int\sqrt{pq}\,\dd u\in[0,1],
\]
the second form because $\int p=\int q=1$. On the quadrature grid with weights $w_k$
the integral is $\tfrac12\sum_kw_k(\sqrt{p_k}-\sqrt{q_k})^2$ after each column has been
normalised under the same weights.
\end{definition}

For a transition kernel every parent state $\gridpt{j}$ gives one density
$\kernel(\cdot\mid\gridpt{j})$, hence one distance $H_j$ between the fitted and the
true kernel. The code reports the median, mean, $90$th percentile and maximum over
parents, the same over the interior parents $|\gridpt{j}|\le\halfwidth/2$ (near the
edge of the domain the true column has mass outside the grid and the distance picks
up truncation), and versions weighted by the parent law, which defaults to the
standard normal because $\Var(a_i)=1$: the weighted mean $\sum_j\pi_jH_j$, the
weighted median, and $\sum_j\pi_jH_j^2$. The second form of the definition is the
one the code evaluates, because $1-\int\sqrt{pq}$ subtracts two numbers near one and
loses everything below $10^{-15}$ to cancellation; the first form returns exactly zero
for identical columns. Which of these summaries the paper's sentence quotes is not
said in the paper; the reader should take the weighted mean as the intended one and
check the generator.

\section{Worked examples}
\label{sec:rev:examples}

\begin{example}[The reverse of a Gaussian, and the bias of a step]
\label{ex:rev:gaussian}
Let $\Pz=\normal01$ in one dimension, so that $\score(x,t)=-x$ at every $t$. The reverse
equation \cref{eq:rev:ou} in the time $\tau$ is
$\dd\bar x=[\bar x-2\bar x]\dd\tau+\sqrt2\,\dd\bar W=-\bar x\,\dd\tau+\sqrt2\,\dd\bar W$:
the forward channel again, for which $\normal01$ is stationary, as it must be. The
Euler step \cref{eq:rev:euler} with constant $h$ is $x\leftarrow(1-h)x+\sqrt{2h}\,\xi$,
and the variance $v$ of the iterate obeys $v\leftarrow(1-h)^2v+2h$, whose fixed point is
\[
  v^\ast=\frac{2h}{1-(1-h)^2}=\frac{2h}{2h-h^2}=\frac{2}{2-h}=1+\tfrac h2+O(h^2).
\]
The scheme's stationary variance is biased by $h/2$: with the first geometric step
$h_0=0.0374$ the bias would be $1.9\%$, with the last, $0.01\%$. This is the weak error
of order one, and it is what the exact arm's log-likelihood gap measures on the chain.
\end{example}

\begin{example}[The residual kurtosis of the frozen chain at $\tmin$]
\label{ex:rev:kurt}
Take $\corr=0.85$, $\ivar=1-\corr^2=0.2775$, a Laplace innovation ($\kappa=3$) and
$\tmin=0.02$, so $e^{-2\tmin}=0.96079$, $e^{-4\tmin}=0.92312$, $\Delta_{\tmin}=0.03921$.
By \cref{prop:rev:kurt},
\[
  \Var(r)=0.96079\times0.2775+0.03921\times1.7225=0.26662+0.06754=0.33416,
\]
\[
  \operatorname{cum}_4(r)=3\times0.92312\times0.2775^2=3\times0.92312\times0.07701=0.21327,
\]
and the excess kurtosis is $0.21327/0.33416^2=0.21327/0.11166=1.91$. A perfect
sampler at $\tmin=0.02$ should therefore show a residual excess kurtosis near $1.9$,
not $3$; the paper's reference row is to be read against this value (see the paper,
Section 5, the generation table). The Gaussian twin shows $0$; a sampler that
reproduces only part of the fourth cumulant lands in between.
\end{example}

\begin{example}[Hellinger distances one can compute by hand]
\label{ex:rev:hellinger}
For $p=\normal01$ and $q=\normal0{\sigma^2}$, $\int\sqrt{pq}=\sqrt{2\sigma/(1+\sigma^2)}$,
so at $\sigma=1.1$: $2.2/2.21=0.99548$, square root $0.99774$, $H^2=0.00226$,
$H=0.048$. For $p=\normal01$ and $q=\normal{\mu}1$, $\int\sqrt{pq}=e^{-\mu^2/8}$, so at
$\mu=0.5$: $H^2=1-e^{-0.03125}=0.0308$, $H=0.175$. A ten per cent error in the width
of an innovation costs less than a half-unit shift of its centre; a column-wise
Hellinger summary is sensitive to both.
\end{example}

\begingroup\setlength{\emergencystretch}{2em}
\where{Generation is Section 5 of the paper (``Generation'') with its table, and the
reverse SDE is named there with its step count, grid and shared Brownian path; the
Hellinger sentence is in Section 4. Experiment E05 (parts \code{generate} and
\code{steps}); the arms come from E02a (fitted kernels) and E03 (network checkpoints).
Code: \code{src/\allowbreak markov\_\allowbreak score/\allowbreak reverse.py} (\code{time\_\allowbreak grid}, \code{reverse\_\allowbreak sde},
\code{reverse\_\allowbreak sde\_\allowbreak with\_\allowbreak increments}, \code{nested\_\allowbreak brownian\_\allowbreak path},
\code{coarsen\_\allowbreak increments}, \code{probability\_\allowbreak flow\_\allowbreak ode}, \code{denoising\_\allowbreak readout}),
\code{src/\allowbreak markov\_\allowbreak score/\allowbreak gen\_\allowbreak metrics.py} (\code{compare} and the functions it calls),
\code{src/\allowbreak markov\_\allowbreak score/\allowbreak metrics.py} (\code{transition\_\allowbreak hellinger}),
\code{experiments/\allowbreak e05\_\allowbreak generation.py}.}
\endgroup

\paragraph{References.}
B.~D.~O.~Anderson, 1982, ``Reverse-time diffusion equation models'', Stochastic
Processes and their Applications 12, 313--326.
U.~G.~Haussmann and E.~Pardoux, 1986, ``Time reversal of diffusions'', Annals of
Probability 14, 1188--1205.
Y.~Song, C.~Durkan, I.~Murray and S.~Ermon, 2021, ``Maximum likelihood training of
score-based diffusion models'', Advances in Neural Information Processing Systems 34.
P.~E.~Kloeden and E.~Platen, 1992, \emph{Numerical Solution of Stochastic Differential
Equations}, Springer.
G.~J.~Sz\'ekely and M.~L.~Rizzo, 2013, ``Energy statistics: a class of statistics based
on distances'', Journal of Statistical Planning and Inference 143, 1249--1272.
J.~Rabin, G.~Peyr\'e, J.~Delon and M.~Bernot, 2012, ``Wasserstein barycenter and its
application to texture mixing'', Scale Space and Variational Methods in Computer
Vision, Lecture Notes in Computer Science 6667, 435--446.
